% =========================================================================
% HƯỚNG DẪN ĐIỀN THÔNG TIN BÁO CÁO NGHIỆM THU ĐỀ TÀI:
% Thay đổi các giá trị trong dấu ngoặc nhọn {...} bên dưới bằng thông tin thật.
% =========================================================================

% Thiết lập trang bìa: điền 1 để chọn mẫu bìa mới (không logo, viền đôi đen sát nhau), 0 để chọn mẫu bìa cũ (có logo, viền xanh)
\BiaMoi{1}



% 1. Thông tin chung về đề tài
\Tendetai{Tên đề tài nghiên cứu khoa học bằng tiếng Việt}
\TendetaiTiengAnh{English Title of the Scientific Research Project}
\MasoDeTai{ABCD-202X-XX}
\ChunhiemDeTai{Học vị.}{Họ và tên chủ nhiệm đề tài}
\newcommand{\KinhphiDuyet}{100} % Số triệu đồng
\newcommand{\KinhphiBangchu}{Một trăm}
\newcommand{\ThoigianHopdong}{12 tháng, từ MM/YYYY đến MM/YYYY}
\newcommand{\ThoigianGiahan}{0 tháng}
\newcommand{\ThoigianThucte}{12 tháng, từ MM/YYYY đến MM/YYYY}
% \ThangNamBaoCao{tháng MM năm YYYY} % Bỏ chọn hoặc để trống để tự động lấy tháng/năm hiện tại
\DiaPhuongBaoCao{Tp. Hồ Chí Minh}

% Chọn loại đề tài (Điền 1 để chọn [x], 0 để bỏ chọn [ ])
\newcommand{\LoaiCoBan}{1}
\newcommand{\LoaiUngDung}{0}

% Chọn ngành/lĩnh vực (Điền 1 để chọn [x], 0 để bỏ chọn [ ])
\newcommand{\NganhToan}{0}
\newcommand{\NganhVatly}{0}
\newcommand{\NganhHoahoc}{0}
\newcommand{\NganhSinhhoc}{0}
\newcommand{\NganhSuckhoe}{0}
\newcommand{\NganhTraidat}{0}
\newcommand{\NganhVatlieu}{0}
\newcommand{\NganhNangluong}{0}
\newcommand{\NganhCokhi}{0}
\newcommand{\NganhDientu}{0}
\newcommand{\NganhCNTT}{1}
\newcommand{\NganhKhac}{0}
\newcommand{\NganhKhacText}{....}

% =========================================================================
% 2. Danh sách chủ nhiệm, thành viên tham gia thực hiện đề tài (Mục 6)
% Cú pháp: \AddMember{Mã số}{Họ và tên}{Đơn vị công tác}{Vai trò trong đề tài}
% Ví dụ: \AddMember{1234}{Nguyễn Văn A}{Khoa CNTT, Trường ĐH KHTN}{Thành viên}
% =========================================================================
\AddMember{1000}{\ChunhiemDeTai}{Trường ĐH KHTN, ĐHQG-HCM}{Chủ nhiệm}
\AddMember{1001}{Học vị. Họ tên thành viên 1}{Trường ĐH KHTN, ĐHQG-HCM}{Thành viên chính}
\AddMember{1002}{Học vị. Họ tên thành viên 2}{Trường ĐH KHTN, ĐHQG-HCM}{Thành viên}

% =========================================================================
% 3. Danh sách công bố khoa học (Phần III, Mục III.1)
% Cú pháp chung cho các loại công bố:
%   - \AddCongBoISI (Tạp chí ISI/quốc tế uy tín)
%   - \AddCongBoIntl (Tạp chí quốc tế khác)
%   - \AddCongBoNat (Tạp chí quốc gia uy tín)
%   - \AddCongBoConfIntl (Kỷ yếu hội nghị quốc tế)
%   - \AddCongBoConfNat (Kỷ yếu hội nghị quốc gia)
%   - \AddCongBoBook (Sách chuyên khảo)
%   - \AddCongBoOther (Kết quả công bố khác)
%
% Cú pháp: \AddCongBo...{Tác giả}{Năm}{Tên công trình}{Thông tin chi tiết}{Xếp hạng}{DOI}{ISSN/ISBN}{Tình trạng}
% Ví dụ:
% \AddCongBoISI{Nguyễn Văn A}{2026}{Tên bài báo}{Tên tạp chí, số, trang}{Q1, IF=5.2}{10.1002/abc}{1234-5678}{Đã công bố}
% =========================================================================
\AddCongBoConfIntl{Họ và tên tác giả 1, Họ và tên tác giả 2}{202X}{Tên công trình công bố bằng tiếng Anh}{Chi tiết về hội nghị khoa học quốc tế, nhà xuất bản, địa điểm tổ chức}{CORE A/B/C hoặc xếp hạng tương đương}{10.1000/xyz123}{XXXX-XXXX}{Đã chấp nhận đăng / Đã xuất bản}

% =========================================================================
% 4. Kết quả đăng ký sở hữu trí tuệ (Phần III, Mục III.2)
% Cú pháp cho các loại sở hữu trí tuệ:
%   - \AddSHTTForeign (Bằng sáng chế nước ngoài)
%   - \AddSHTTDomestic (Bằng sáng chế trong nước)
%   - \AddSHTTUtility (Giải pháp hữu ích)
%   - \AddSHTTSoftware (Bản quyền phần mềm)
%   - \AddSHTTOther (Các kết quả khác)
%
% Cú pháp: \AddSHTT...{Tác giả}{Tên công trình}{Nơi nhận hồ sơ/Cấp chứng nhận, số và ngày quyết định}{Tình trạng}
% Ví dụ:
% \AddSHTTSoftware{Nguyễn Văn A}{Phần mềm quản lý}{Cục Bản quyền Tác giả, QĐ số 456 ngày 01/01/2026}{Đã cấp}
% =========================================================================
% \AddSHTTSoftware{Tác giả 1, Tác giả 2}{Tên phần mềm/sáng chế}{Cục Bản quyền tác giả, QĐ số...}{Đã cấp}

% =========================================================================
% 5. Kết quả đào tạo (Phần III, Mục III.3)
% Cú pháp cho các loại đào tạo:
%   - \AddDaoTaoPHD (Nghiên cứu sinh)
%   - \AddDaoTaoMaster (Học viên cao học)
%   - \AddDaoTaoUndergrad (Sinh viên đại học)
%
% Cú pháp: \AddDaoTao...{Họ và tên}{Thời gian tham gia (từ mm/yyyy đến mm/yyyy)}{Tên đề tài luận án/luận văn/khóa luận}{Ghi chú}
% Ví dụ:
% \AddDaoTaoUndergrad{Nguyễn Văn B}{từ 11/2025 đến 06/2026}{Nghiên cứu về Hate Speech}{Sinh viên làm khóa luận tốt nghiệp}
% =========================================================================
% \AddDaoTaoMaster{Họ và tên học viên}{từ MM/YYYY đến MM/YYYY}{Tên đề tài luận văn thạc sĩ}{Ghi chú đào tạo}

% =========================================================================
% 6. Tình hình sử dụng kinh phí (Phần IV)
% Định dạng: \SetKinhPhi...{Kinh phí được duyệt}{Kinh phí thực hiện}
% =========================================================================
\SetKinhPhiNhanCong{80000000}{80000000}
\SetKinhPhiVatLieu{10000000}{10000000}
\SetKinhPhiThietBi{5000000}{5000000}
\SetKinhPhiDiLai{3000000}{3000000}
\SetKinhPhiThueNgoai{0}{0}
\SetKinhPhiTrucTiepKhac{0}{0}
\SetKinhPhiQuanLy{2000000}{2000000}{Quản lý chung}
\SetKinhPhiHoiDong{0}{0}{}
